\documentclass[11pt,a4paper]{article}
\usepackage{vespa_style}

\begin{document}

\vespatitlepage
  {Bibliographie}
  {Sources, normes et documentation technique}
  {Projet \vespa{} mk.2 — ENSIL-ENSCI, Limoges}

\tableofcontents
\newpage

% =====================================================================
\section{Organisation et usage}

Cette bibliographie n'est pas une liste de lectures : c'est l'ensemble des documents dont le
projet \textbf{dépend effectivement}. Chaque entrée est accompagnée, dans le fichier source,
d'une note précisant \textbf{à quoi elle a servi} — quel calcul elle fonde, quel composant
elle décrit, quelle contrainte elle impose. Une référence qui n'a rien déterminé n'y figure
pas.

\paragraph{Trois natures de sources, trois usages.}

\begin{enumerate}
  \item \textbf{Les articles scientifiques} (\fichier{docs/references/articles/}) fondent le
        \emph{cahier des charges} : biologie et comportement de la cible, état de l'art de la
        détection. Ils déterminent \emph{ce que le système doit faire}, et surtout ce qu'il
        \emph{ne doit pas} faire — ne pas tirer sur \emph{Vespa crabro}.
  \item \textbf{Les manuels constructeur et fiches techniques}
        (\fichier{docs/references/datasheets/}) fondent la \emph{conception matérielle}. Ce
        sont eux qui fixent les brochages, les seuils électriques et les limites physiques.
        Plusieurs verrous du projet sont des lignes de ces documents qui n'avaient pas été
        lues assez tôt.
  \item \textbf{Les normes} (ISO 11898 pour le bus CAN, IEC 60825 pour le laser) sont
        \emph{prescriptives}. Celle de sécurité laser, en particulier, ne se discute pas.
\end{enumerate}

\begin{note}[Vérification des métadonnées]
Les métadonnées de chaque article (auteurs, revue, volume, DOI) ont été relevées
\textbf{sur la première page du document lui-même}, et non sur une base de données
secondaire. Les fiches techniques renvoient au matériel identifié sans ambiguïté par
\fichier{platformio.ini}, \fichier{pinout.h} et la nomenclature.
\end{note}

\paragraph{Fichier source.} L'ensemble est maintenu en BibTeX dans
\textbf{\fichier{docs/src/vespa.bib}}. C'est ce fichier qu'il faut enrichir, et non ce PDF,
qui en est une vue commentée.

% =====================================================================
\section{Biologie de la cible}

Le projet vise une espèce précise, et cette précision est une \textbf{exigence éthique}, pas
un détail : il s'agit de neutraliser \emph{Vespa velutina nigrithorax} (invasive) sans
toucher \emph{Vespa crabro} (le frelon européen, natif et protégé).

\begin{itemize}[leftmargin=*]
  \item \textbf{O'Shea-Wheller \emph{et al.} (2024)} \cite{osheawheller2024vespai} —
        \textsc{VespAI}. C'est \emph{la} référence structurante du projet : le modèle de
        détection initialement retenu (YOLOv5s + ResNet, précision-rappel $\geq 0{,}99$).
        C'est aussi la source du principal verrou : \textbf{il est entraîné en RGB}, et nos
        capteurs sont monochromes.

  \item \textbf{Sipos \emph{et al.} (2025)} \cite{sipos2025morphologie} — morphologie
        comparée \emph{velutina} / \emph{crabro} par micro-CT. Double usage : les
        performances de vol (vitesse, manœuvrabilité, vol stationnaire) alimentent le
        dimensionnement cinématique ; et l'article établit que la discrimination entre les
        deux espèces repose largement sur la \textbf{coloration} — ce qui explique pourquoi
        elle est structurellement compromise en monochrome.

  \item \textbf{Perrard \emph{et al.} (2009)} \cite{perrard2009colonie} — activité de la
        colonie. Documente le comportement de \textbf{vol de guet} devant les ruches.
        C'est l'hypothèse opérationnelle qui rend le système viable à 60~fps : on ne traque
        pas un frelon en vol de croisière, on traque un chasseur en vol quasi stationnaire.

  \item \textbf{Lioy \emph{et al.} (2023)} \cite{lioy2023niche} — recouvrement de niche
        climatique avec les espèces natives. Fonde le contexte d'invasion et l'exigence de
        sélectivité.

  \item \textbf{Villemant \emph{et al.} (2015)} \cite{villemant2015parasites} — lutte
        biologique par nématodes mermithidés. Contrepoint utile : aucune méthode de
        régulation existante n'est à la fois sélective et efficace, ce qui justifie
        d'explorer la lutte active.

  \item \textbf{Ruiz-Cristi \emph{et al.} (2020)} \cite{ruizcristi2020thermotolerance} —
        \textbf{tolérance thermique de \emph{V. velutina nigrithorax}}. Mesurée sur
        l'espèce cible elle-même, tous sexes, castes et stades confondus : la limite
        thermique létale se situe autour de \textbf{45\,\textcelsius{}}. \textbf{C'est la
        donnée biologique sur laquelle repose le dimensionnement de l'effecteur} — la
        température que le tir doit atteindre. Corroborée par
        \textbf{Ono \emph{et al.} (1995)} \cite{ono1995thermal}, qui montre que l'abeille
        japonaise exploite précisément cet écart de tolérance (frelon
        $\approx 45$\,\textcelsius{}, abeille $\approx 50$\,\textcelsius{}) pour tuer le
        frelon par hyperthermie dans une « boule » à 47\,\textcelsius{}. La neutralisation
        thermique du frelon n'est donc pas une extrapolation : c'est un mécanisme attesté
        dans la nature.
\end{itemize}

\begin{note}[Ce que ces deux articles ne disent pas]
Ils donnent un \textbf{seuil de température}, pas une \textbf{dose optique}. Leurs temps
d'exposition (13~s à plusieurs minutes) relèvent d'un chauffage \emph{convectif du corps
entier} — étuve, flux d'air humide. Ils ne se transposent pas tels quels à une impulsion
laser localisée de 20~ms. Le passage du seuil à la dose est l'objet des deux références
\emph{Photonic Fence} du chapitre « Optique et sécurité laser », et ce passage
\textbf{reste incomplet} : voir l'avertissement qui s'y trouve.
\end{note}

% =====================================================================
\section{Vision par ordinateur}

\begin{itemize}[leftmargin=*]
  \item \textbf{Garrido-Jurado \emph{et al.} (2014)} \cite{garridojurado2014aruco} —
        marqueurs \textsc{ArUco}. \textbf{Solution effectivement déployée} en remplacement
        de VespAI : haut contraste, nativement adaptée aux capteurs monochromes, aucun
        entraînement, pose 3D exacte. Complété par
        Romero-Ramirez \emph{et al.} (2018) \cite{romeroramirez2018speeded} pour la
        détection accélérée (algorithme du module \texttt{cv::aruco}).

  \item \textbf{Zhang (2000)} \cite{zhang2000calibration} — calibration de caméra.
        Méthode implémentée par \texttt{cv::calibrateCamera} / \texttt{cv::stereoCalibrate},
        utilisée par l'outil \fichier{software/tools/calib\_stereo}.

  \item \textbf{Hartley \& Zisserman (2004)} \cite{hartley2004multiview} — géométrie
        épipolaire et triangulation. Référence de \texttt{StereoTriangulator}.

  \item \textbf{Kalman (1960)} \cite{kalman1960filter} — filtrage prédictif. Fondement de
        \texttt{PredictiveTurretSight}, qui compense la latence propre de la chaîne
        vision $\rightarrow$ CAN $\rightarrow$ asservissement. Sans cette prédiction, la
        tourelle vise systématiquement là où la cible \emph{était}.

  \item \textbf{Ultralytics YOLOv5} \cite{ultralytics2020yolov5} — architecture
        sous-jacente de VespAI. Le dépôt est \emph{requis} pour désérialiser les poids
        \fichier{.pt} (qui référencent le module Python \texttt{models}) — détail non
        évident qui a coûté du temps.

  \item \textbf{OpenCV} \cite{opencv}, \textbf{V4L2} \cite{v4l2},
        \textbf{TensorRT} \cite{nvidia_tensorrt} — bibliothèques et interfaces
        effectivement employées. V4L2 (\texttt{mmap}) porte la chaîne d'acquisition
        « zéro-copie ».
\end{itemize}

% =====================================================================
\newpage
\section{Commande moteur}

\begin{itemize}[leftmargin=*]
  \item \textbf{SimpleFOC} \cite{simplefoc} — bibliothèque d'asservissement vectoriel
        (version 2.3.2) utilisée par le \fichier{motion\_node}.

  \item \textbf{ST UM2538} \cite{st_um2538_focsdk} — théorie de la commande vectorielle
        (transformées de Clarke et Park) appliquée au matériel ST.

  \item \textbf{ST UM2415} \cite{st_um2415_ihm16m1} — étage de puissance
        X-NUCLEO-IHM16M1 (STSPIN830). Deux exemplaires, cavaliers en FOC 3 shunts.

  \item \textbf{ST UM2505} \cite{st_um2505_nucleog431} — carte NUCLEO-G431RB.
        \textbf{C'est ce document qui établit que \broche{PA2} est câblée sur la liaison
        série virtuelle du ST-LINK} — origine d'un verrou coûteux (voir « Difficultés
        techniques »).

  \item \textbf{ams AS5048A} \cite{ams_as5048a} — encodeur magnétique absolu 14 bits.
        Fixe la borne théorique de précision : $360/2^{14} \approx 0{,}022^\circ$.
\end{itemize}

% =====================================================================
\section{Capteurs et calculateur}

\begin{itemize}[leftmargin=*]
  \item \textbf{OmniVision OV9281} \cite{omnivision_ov9281} — capteur monochrome à
        obturateur global $1\,280 \times 800$. Le choix du monochrome n'est pas une
        préférence : \textbf{aucun capteur trichrome à obturateur global n'existe à prix
        abordable}. C'est la cause racine de l'incompatibilité avec VespAI.

  \item \textbf{Arducam UC-788} \cite{arducam_uc788} — carte porteuse assurant la
        compatibilité avec l'écosystème Jetson et le déclenchement externe. \textbf{C'est la
        source qui établit que le constructeur ne prévoit pas la synchronisation matérielle
        sur cartes standard} : elle passe par son accessoire propriétaire (\emph{Camarray
        HAT}), et la piste \texttt{XVS/FSIN} est coupée en sortie d'usine sur les cartes
        vendues seules. Le projet obtient malgré tout un déclenchement stéréo matériel sur
        deux cartes standard (voir « Difficultés techniques »).

  \item \textbf{NVIDIA Jetson Orin Nano — carte porteuse} \cite{nvidia_orin_nano_carrier} —
        brochage de l'en-tête 40 points. Source des broches PWM (32/33) utilisées pour la
        synchronisation stéréo, \textbf{et} des broches CAN0 (\broche{PAA0}/\broche{PAA1},
        en-tête 31/29) sur lesquelles le bus a fonctionné. \textbf{C'est ce document qui
        établit que CAN0 sort sur les broches 29/31, et non celles de CAN1} — la confusion
        entre les deux a coûté plusieurs itérations d'\emph{overlay}.

  \item \textbf{JetPack / Jetson Linux} \cite{nvidia_jetpack} — environnement déployé
        (JetPack 6.2). \textbf{Le noyau doit être remplacé par le noyau Arducam.}
\end{itemize}

% =====================================================================
\section{Communication}

\begin{itemize}[leftmargin=*]
  \item \textbf{ISO 11898-1:2015} \cite{iso11898_1} — couche liaison CAN. L'arbitrage bit à
        bit non destructif fonde toute la hiérarchie de priorités du protocole VESPA :
        \texttt{FIRE\_CMD} porte l'identifiant le plus faible (\texttt{0x010}), donc la
        priorité absolue sur le bus.

  \item \textbf{SocketCAN} \cite{socketcan} — pile CAN de Linux, côté Jetson. Son interface
        virtuelle \texttt{vcan0} a permis de développer et de valider tout le protocole
        \textbf{avant} que le bus physique soit disponible ; le code applicatif est passé sur
        \texttt{can0} réel \textbf{sans une ligne de changement}. Ses compteurs
        (\texttt{ip -s link show can0}) sont aussi l'instrument qui a mis en évidence
        l'asymétrie réception/émission révélatrice des transceivers contrefaits.
\end{itemize}

% =====================================================================
\section{Optique et sécurité laser}

\begin{itemize}[leftmargin=*]
  \item \textbf{IEC 60825-1:2014} \cite{iec60825} — classification des produits laser.
        \textbf{La norme qui commande tout le dispositif de sécurité du projet.}
        L'effecteur (diode bleue multiwatt) relève de la \textbf{Classe 4} : dangereux
        y compris en réflexion diffuse, risque d'incendie, aucune distance sûre en
        intérieur. C'est la raison pour laquelle la diode n'a jamais été mise sous tension
        en l'absence de watchdog.

  \item \textbf{Hecht, \emph{Optics}} \cite{hecht2016optics} — relations de conjugaison de
        Descartes, utilisées pour établir la loi de commande du mécanisme de collimation
        dynamique (deux lentilles minces convergentes).

  \item \textbf{Siegman, \emph{Lasers}} \cite{siegman1986lasers} — divergence et irradiance
        du faisceau. Support du tracé de densité de puissance le long de l'axe optique
        (\fichier{simulation/optique/Laser\_Optics\_v1.m}).

  \item \textbf{Keller \emph{et al.} (2020)} \cite{keller2020photonicfence} —
        \textsc{Photonic Fence}. \textbf{Le seul système publié qui neutralise au laser des
        insectes \emph{en vol}}, avec des courbes dose-réponse mesurées. Il fonde
        \textbf{deux} des trois constantes de tir de
        \fichier{simulation/optique/Laser\_Optics\_v1.m} :
        \begin{itemize}
          \item \textbf{le choix du bleu} — les longueurs d'onde visibles exigent une
                exposition bien plus faible que le proche infrarouge : LD90
                $= 1{,}6$~J/cm$^2$ à 445~nm, contre $12{,}9$~J/cm$^2$ à 1064~nm. Le
                « couplage optimal avec la cuticule » invoqué par l'annexe sécurité a donc
                une base mesurée ;
          \item \textbf{la durée d'impulsion} — 25~ms est identifiée comme la durée idéale
                pour minimiser la puissance requise, et le comportement de vol de la cible
                n'affecte pas la mortalité en deçà. Les \textbf{20~ms} de
                \texttt{T\_exposure} tombent dans cette fenêtre.
        \end{itemize}
        Complété par \textbf{Keller \emph{et al.} (2016)} \cite{keller2016mosquito}, l'étude
        de dose antérieure (sujets anesthésiés, balayage 532--10\,600~nm, 10~ns--25~ms) :
        c'est la source de la \emph{méthode} — comment on \emph{établit} une fluence létale
        au lieu de la postuler.
\end{itemize}

\begin{avertissement}[La fluence létale de 50~J/cm$^2$ n'est toujours pas sourcée]
Les deux références ci-dessus fondent la longueur d'onde et la durée d'impulsion. Elles ne
fondent \textbf{pas} la troisième constante, et c'est la plus lourde de conséquences :
\texttt{Fluence\_Lethal = 50 J/cm\textasciicircum 2} dans
\fichier{simulation/optique/Laser\_Optics\_v1.m}.

\begin{itemize}
  \item C'est \textbf{environ 30 fois} la LD90 publiée dans le bleu ($1{,}6$~J/cm$^2$) —
        mais celle-ci est mesurée sur un \emph{Anopheles}, deux ordres de grandeur plus
        léger qu'une ouvrière de \emph{V. velutina}.
  \item \textbf{Aucune mesure de létalité laser n'existe pour un frelon.} Keller \emph{et
        al.} ne testent que moustiques, drosophiles et psylles : rien qui approche la masse
        de la cible du projet. L'extrapolation est donc inévitable — mais elle doit être
        \textbf{déclarée comme telle}.
  \item Le sens de l'erreur n'est pas neutre. Surestimer la fluence létale est prudent pour
        l'\emph{efficacité}, et \textbf{dangereux pour la sécurité} : c'est ce chiffre qui
        commande la puissance requise, donc l'irradiance de $\approx 2\,500$~W/cm$^2$, donc
        la Classe 4 et les \textbf{200~m} de portée du danger oculaire. Une valeur plus
        juste rendrait peut-être le dispositif moins dangereux.
\end{itemize}

\textbf{Travail à reprendre :} établir la dose létale pour \emph{V. velutina} en partant du
seuil thermique de Ruiz-Cristi \emph{et al.} \cite{ruizcristi2020thermotolerance}
($\approx 45$\,\textcelsius{}) et des paramètres morphologiques de Sipos \emph{et al.}
\cite{sipos2025morphologie} (masse et surface du thorax), par la méthode de
Keller \emph{et al.} (2016) \cite{keller2016mosquito}. Tant que ce calcul n'est pas refait,
\textbf{les 50~J/cm$^2$ restent une hypothèse, pas un résultat.}
\end{avertissement}

% =====================================================================
\section{Travaux antérieurs}

\begin{itemize}[leftmargin=*]
  \item \textbf{Cosson \& Grimal (2025)} \cite{cosson2025vespa_mk1} — \textbf{prototype
        mk.1}, « Dispositif de lutte contre les frelons asiatiques — Système de pointage de
        la cible ». Détection par YOLOv8, tourelle à moteurs pas-à-pas pilotée par Arduino
        Uno (shield DRV8825), liaison Python/Arduino par port série, et une première
        exploration du portage sur GPU Jetson.

        \textbf{C'est la référence fondatrice du mk.2}, et elle se lit à l'envers : les
        limites qu'elle constate définissent en creux notre cahier des charges. L'absence de
        FPU sur l'ATmega328P et le décrochage des pas-à-pas ont donné la commande vectorielle
        sur encodeurs absolus ; l'absence de perception de la profondeur a donné la
        stéréo-vision et la triangulation. \textbf{Le mk.2 n'est pas une amélioration du
        mk.1 : c'est la réponse à son bilan.}

        Archivé dans
        \fichier{docs/references/rapports/Cosson\_Grimal\_2025\_VESPA\_mk1\_Pointage.pdf}.

  \item \textbf{Delmas, Desmedt \& Gilson (2026)} \cite{delmas2026vespa_mk2} — le rapport de
        soutenance du présent prototype, conservé tel quel dans
        \fichier{docs/Rapport\_PFE\_2026\_VESPA.pdf}.
\end{itemize}

\begin{note}[Filiation des deux prototypes]
Les deux rapports se lisent ensemble. Le mk.1 établit la \emph{faisabilité} de la détection
et du pointage ; le mk.2 attaque les trois limites qu'il laisse ouvertes — la précision de
l'asservissement, la mesure de la distance, et la puissance de l'effecteur. Un repreneur qui
n'aurait à lire qu'un seul document d'histoire du projet doit lire le mk.1 en premier.
\end{note}

% =====================================================================
\newpage
\section*{Références complètes}
\addcontentsline{toc}{section}{Références complètes}

\nocite{*}
\bibliographystyle{plain}
\bibliography{vespa}

\end{document}
